\clearpage\pagestyle{otMain}\setlength{\otBlocEspacement}{2.5mm}\fontsize{11}{15}\selectfont\corrigesChapitre{14}{Médiatrice et équidistance}
\begin{blocLivre}{cor-14-1}\begin{corrige}{1}{Reconnaître la médiatrice}\demarche{Vérifier les deux conditions de la définition.}
$(d_1)$ est perpendiculaire mais ne passe pas par le milieu ; $(d_2)$ passe par le milieu mais n’est pas perpendiculaire ; $(d_4)$ ne vérifie aucune condition. La médiatrice est $\res{(d_3)}$ : perpendiculaire en son milieu.
\label{fig-cor-14-1}\figureGeo[corrige]{med-1}
\end{corrige}\end{blocLivre}
\begin{blocLivre}{cor-14-2}\begin{corrige}{2}{Coder une médiatrice}\demarche{Coder le milieu et l’angle droit.}
On code les deux moitiés de [AB] par un même signe et l’angle droit entre (d) et (AB) par un petit carré, en vert.
\label{fig-cor-14-2}\figureGeo[corrige]{med-2}
\end{corrige}\end{blocLivre}
\begin{blocLivre}{cor-14-3}\begin{corrige}{3}{Construire à l’équerre}\demarche{Milieu d’abord, perpendiculaire ensuite.}
On place le milieu I avec $AI=3$ cm, puis on trace la perpendiculaire à (AB) passant par I : c’est la médiatrice, en vert.
\label{fig-cor-14-3}\figureGeo[corrige]{med-3}
\end{corrige}\end{blocLivre}
\begin{blocLivre}{cor-14-4}\begin{corrige}{4}{Une longueur sans mesurer}\demarche{Utiliser la propriété dans le sens direct.}
M est sur la médiatrice de [AB], donc il est à égale distance de A et de B : $\res{MB=4{,}2\ \mathrm{cm}}$.
\label{fig-cor-14-4}\figureGeo[corrige]{med-4}
\end{corrige}\end{blocLivre}
\begin{blocLivre}{cor-14-5}\begin{corrige}{5}{Sur la médiatrice ?}\demarche{Utiliser la propriété réciproque.}
P est à égale distance de A et de B : il est $\res{\text{sur la médiatrice}}$. Comme $QA\ne QB$, Q n’est $\res{\text{pas sur la médiatrice}}$.
\label{fig-cor-14-5}\figureGeo[corrige]{med-5}
\end{corrige}\end{blocLivre}
\begin{blocLivre}{cor-14-6}\begin{corrige}{6}{Construire au compas}\demarche{L’écartement doit dépasser la moitié de CD.}
Les arcs de rayon 4 cm, de centres C et D, se coupent en deux points E et F, à 4 cm de C et de D. La droite (EF) est la médiatrice de [CD], en vert.
\label{fig-cor-14-6}\figureGeo[corrige]{med-6}
\end{corrige}\end{blocLivre}
\begin{blocLivre}{cor-14-7}\begin{corrige}{7}{Le triangle isocèle}\demarche{Rédiger avec la propriété réciproque.}
Je sais que $AB=AC$. Or un point à égale distance de B et de C est sur la médiatrice de [BC]. Donc $\res{\text{A est sur la médiatrice de [BC]}}$. Elle coupe [BC] en son milieu, à 3 cm de B.
\label{fig-cor-14-7}\figureGeo[corrige]{med-7}
\end{corrige}\end{blocLivre}
\begin{blocLivre}{cor-14-8}\begin{corrige}{8}{Le losange}\demarche{Trouver deux points de la médiatrice.}
$BA=BC$, donc B est sur la médiatrice de [AC]. $DA=DC$, donc D est aussi sur cette médiatrice. Une droite est déterminée par deux de ses points : la médiatrice de [AC] est $\res{\text{la droite (BD)}}$.
\label{fig-cor-14-8}\figureGeo[corrige]{med-8}
\end{corrige}\end{blocLivre}
\begin{blocLivre}{cor-14-9}\begin{corrige}{9}{Un point sur une droite}\demarche{Remplacer la condition de distance par une droite.}
M est à égale distance de A et de B, donc il est sur la médiatrice de [AB]. On trace cette médiatrice : elle coupe (d) au point M cherché, en vert.
\label{fig-cor-14-9}\figureGeo[corrige]{med-9}
\end{corrige}\end{blocLivre}
\begin{blocLivre}{cor-14-10}\begin{corrige}{10}{Des cercles passant par A et B}\demarche{Un cercle passant par A et B a deux rayons égaux, OA et OB.}
Le centre O d’un tel cercle vérifie $OA=OB$ (deux rayons) : il est sur la $\res{\text{médiatrice de [AB]}}$. Deux exemples sont tracés en vert, de centres $O_1$ et $O_2$.
\label{fig-cor-14-10}\figureGeo[corrige]{med-10}
\end{corrige}\end{blocLivre}
\begin{blocLivre}{cor-14-11}\begin{corrige}{11}{L’arrêt de bus}\demarche{Traduire l’équidistance par la médiatrice.}
S doit vérifier $SA=SB$ : il est sur la médiatrice de [AB]. Le point S cherché est l’intersection de cette médiatrice et de la route, en vert.
\label{fig-cor-14-11}\figureGeo[corrige]{med-11}
\end{corrige}\end{blocLivre}
\begin{blocLivre}{cor-14-12}\begin{corrige}{12}{Plus près de A ou de B ?}\demarche{La médiatrice partage le plan en deux zones.}
La frontière est formée des points à égale distance de A et de B : c’est la $\res{\text{médiatrice de [AB]}}$. Les points plus proches de A sont du côté de A, colorés en vert.
\label{fig-cor-14-12}\figureGeo[corrige]{med-12}
\end{corrige}\end{blocLivre}
\begin{blocLivre}{cor-14-13}\begin{corrige}{13}{Le pliage}\demarche{Relier pliage et équidistance.}
Chaque point du pli est à la même distance de A et de B, puisque le pliage superpose A et B. Le pli est donc formé de points équidistants de A et de B : c’est la $\res{\text{médiatrice de [AB]}}$, en vert.
\label{fig-cor-14-13}\figureGeo[corrige]{med-13}
\end{corrige}\end{blocLivre}
\begin{blocLivre}{cor-14-14}\begin{corrige}{14}{Raisonner}\demarche{Choisir le bon sens de la propriété.}
C est sur la médiatrice, donc $\res{CB=CA=5\ \mathrm{cm}}$. D est à égale distance de A et de B, donc $\res{\text{D est sur (d)}}$. La première réponse utilise la propriété, la seconde sa réciproque.
\label{fig-cor-14-14}\figureGeo[corrige]{med-14}
\end{corrige}\end{blocLivre}
\begin{blocLivre}{cor-14-15}\begin{corrige}{15}{Vrai ou faux}\demarche{Penser à tous les points de la médiatrice.}
$\res{\text{Non}}$ : tous les points de la médiatrice sont à égale distance de A et de B, comme le point M de la figure. Un seul d’entre eux est sur le segment [AB] : son milieu I.
\label{fig-cor-14-15}\figureGeo[corrige]{med-15}
\end{corrige}\end{blocLivre}
